\begin{problem}[A probability table with unknown constant]
A discrete random variable $X$ has probability distribution
\[
P(X=0)=k,\qquad P(X=1)=2k,\qquad P(X=2)=3k,\qquad P(X=3)=4k.
\]
\begin{enumerate}[(a)]
  \item Find $k$.
  \item Find $E(X)$.
  \item Find $\mathrm{Var}(X)$.
\end{enumerate}
\end{problem}

\begin{solution}
\textbf{(a)} The probabilities must sum to $1$, so
\[
k+2k+3k+4k=10k=1.
\]
Hence
\[
k=\frac1{10}.
\]

\textbf{(b)} Then
\[
P(X=0)=\frac1{10},\quad
P(X=1)=\frac2{10},\quad
P(X=2)=\frac3{10},\quad
P(X=3)=\frac4{10}.
\]

Now
\[
E(X)=0\cdot \frac1{10}+1\cdot \frac2{10}+2\cdot \frac3{10}+3\cdot \frac4{10}
=\frac{0+2+6+12}{10}=2.
\]

\textbf{(c)} Next,
\[
E(X^2)=0^2\cdot \frac1{10}+1^2\cdot \frac2{10}+2^2\cdot \frac3{10}+3^2\cdot \frac4{10}
=\frac{0+2+12+36}{10}=5.
\]
Therefore,
\[
\mathrm{Var}(X)=E(X^2)-[E(X)]^2=5-2^2=1.
\]
\end{solution}

\begin{takeaways}
\item A probability table must add to $1$ before anything else.
\item Variance is often quickest via $E(X^2)-[E(X)]^2$.
\item Keep expectation and variance conceptually separate: centre versus spread.
\end{takeaways}

\begin{problem}[Exactly two successes in a binomial setting]
A biased coin lands heads with probability $0.7$. The coin is tossed $5$ times. Find the probability of getting exactly $2$ heads.
\end{problem}

\begin{solution}
Let $X$ be the number of heads. Then
\[
X\sim \mathrm{Bin}(5,0.7).
\]

We want
\[
P(X=2)=\binom52(0.7)^2(0.3)^3.
\]

Since
\[
\binom52=10,
\]
we obtain
\[
P(X=2)=10(0.49)(0.027)=0.1323.
\]

Hence the probability is
\[
\boxed{0.1323}.
\]
\end{solution}

\begin{takeaways}
\item Check the four binomial conditions before applying the formula.
\item The combination counts where the successes are placed.
\item Keep the success probability attached to the event being counted.
\end{takeaways}

\begin{problem}[At least one defective item]
Each item produced by a machine is defective with probability $0.04$, independently of other items. A sample of $8$ items is chosen.
\begin{enumerate}[(a)]
  \item Find the probability that exactly one item is defective.
  \item Find the probability that at least one item is defective.
\end{enumerate}
\end{problem}

\begin{solution}
Let $X$ be the number of defective items. Then
\[
X\sim \mathrm{Bin}(8,0.04).
\]

\textbf{(a)} For exactly one defective item,
\[
P(X=1)=\binom81(0.04)^1(0.96)^7.
\]
So
\[
P(X=1)=8(0.04)(0.96)^7\approx 0.2409.
\]

\textbf{(b)} For at least one defective item, use the complement:
\[
P(X\ge 1)=1-P(X=0)=1-(0.96)^8\approx 0.2786.
\]
\end{solution}

\begin{takeaways}
\item In binomial questions, ``at least one'' is usually a complement.
\item Exact values and cumulative values often call for different tactics.
\item Binomial modelling becomes natural when defect decisions are repeated independently.
\end{takeaways}

\begin{problem}[A binomial distribution from a target score]
In a multiple-choice quiz, each question has four options and exactly one correct answer. A student guesses every answer independently on $6$ questions.
\begin{enumerate}[(a)]
  \item Find the probability that the student answers exactly $4$ questions correctly.
  \item Find the probability that the student answers at least $1$ question correctly.
\end{enumerate}
\end{problem}

\begin{solution}
Let $X$ be the number of correct answers. Since each question is guessed independently and
\[
P(\text{correct})=\frac14,
\]
we have
\[
X\sim \mathrm{Bin}\left(6,\frac14\right).
\]

\textbf{(a)} Then
\[
P(X=4)=\binom64\left(\frac14\right)^4\left(\frac34\right)^2.
\]
Since
\[
\binom64=15,
\]
we get
\[
P(X=4)=15\cdot \frac1{256}\cdot \frac9{16}
=\frac{135}{4096}.
\]

\textbf{(b)} Using a complement,
\[
P(X\ge 1)=1-P(X=0)=1-\left(\frac34\right)^6.
\]
Hence
\[
P(X\ge 1)=1-\frac{729}{4096}=\frac{3367}{4096}.
\]
\end{solution}

\begin{takeaways}
\item Guessing questions independently creates a natural binomial model.
\item Exact counts use the binomial formula directly.
\item ``At least one'' remains a complement even in a binomial setting.
\end{takeaways}

\begin{problem}[Comparing two games by expectation]
Game A pays \$8 with probability $0.2$ and \$0 otherwise. Game B pays according to:
\[
P(Y=1)=0.3,\qquad P(Y=4)=0.4,\qquad P(Y=7)=0.3.
\]
Each game costs \$2 to play once. Determine which game has the higher expected profit.
\end{problem}

\begin{solution}
Let $A$ be the payout from Game A. Then
\[
E(A)=8(0.2)+0(0.8)=1.6.
\]
So the expected profit from Game A is
\[
1.6-2=-0.4.
\]

Now let $Y$ be the payout from Game B. Then
\[
E(Y)=1(0.3)+4(0.4)+7(0.3)=0.3+1.6+2.1=4.
\]
So the expected profit from Game B is
\[
4-2=2.
\]

Therefore Game B has the higher expected profit.
\end{solution}

\begin{takeaways}
\item Separate payout from profit when a game has an entry cost.
\item Expected value gives a fair comparison over many plays.
\item A game can have an attractive top prize and still have poor expected profit.
\end{takeaways}
